\documentclass[12pt,a4paper]{exampaper}
\examname{概率论与数理统计}

\begin{document}

\examtitle{东南大学试卷（A 卷）}

\vspace{0.8em}

\begin{center}
  \examfield[3.4cm]{课程名称}{概率论与数理统计}\hspace{0.4em}%
  \examfield[2.2cm]{课程代码}{ST203}\hspace{0.4em}%
  \examfield[1.8cm]{考试学期}{2025-2026-2}\\[0.4em]
  \examfield[4.4cm]{适用专业}{信息工程}\hspace{0.4em}%
  \examfield[2.0cm]{考试形式}{闭\hspace{0.5em}卷}\hspace{0.4em}%
  \examfield[1.9cm]{考试时长}{120\hspace{0.4em}分钟}
\end{center}

\vspace{0.8em}

\examscoretable[4]

\vspace{0.6em}

\examsection{单项选择题（本题共 2 小题，每小题 4 分，满分 8 分）}
\begin{examquestions}
  \examquestion 设 $X \sim N(\mu, \sigma^2)$，则 $P(|X-\mu| < 2\sigma)$ 约为 \quad A.\ $0.6826$ \quad B.\ $0.9544$ \quad C.\ $0.9974$ \quad D.\ $0.5000$
  \examquestion 设 $A,B$ 为随机事件，且 $P(A)=0.5$，$P(B)=0.4$，$P(AB)=0.2$，则 $P(A \cup B)=$ \quad A.\ $0.5$ \quad B.\ $0.6$ \quad C.\ $0.7$ \quad D.\ $0.9$
\end{examquestions}

\vspace{0.6em}

\examsection{填空题（本题共 2 小题，每小题 4 分，满分 8 分）}
\begin{examquestions}
  \examquestion 设总体 $X$ 的方差为 $\sigma^2$，样本容量为 $n$，则样本均值 $\bar{X}$ 的方差为 $\underline{\hspace{2.5cm}}$。
  \examquestion 若 $X$ 服从参数为 $\lambda$ 的泊松分布，则 $E(X)=$ $\underline{\hspace{2.5cm}}$。
\end{examquestions}

\vspace{0.6em}

\examsection{计算题（本题共 2 小题，每小题 12 分，满分 24 分）}
\begin{examquestions}[0pt]
  \examquestion 设总体 $X \sim N(\mu, \sigma^2)$，$\mu$ 未知，$\sigma^2$ 已知。抽取容量为 $n$ 的样本，样本均值为 $\bar{x}$。试在显著性水平 $\alpha=0.05$ 下检验假设 $H_0: \mu = \mu_0$ 与 $H_1: \mu \neq \mu_0$，并写出拒绝域。
  \par\vspace{7cm}
  \examquestion 设随机变量 $X$ 的概率密度函数为 $f(x) = \begin{cases} \lambda e^{-\lambda x}, & x > 0 \\ 0, & x \leq 0 \end{cases}$，求 $E(X)$ 与 $\mathrm{Var}(X)$。
  \par\vspace{7cm}
\end{examquestions}

\vspace{0.6em}

\examsection{简答题（本题共 1 小题，每小题 10 分，满分 10 分）}
\begin{examquestions}[0pt]
  \examquestion 叙述中心极限定理（林德伯格-列维形式），并说明它在统计推断中的作用。
  \par\vspace{10cm}
\end{examquestions}

\end{document}
